\chapter{Pipeline Python : Le DataPreprocessor}
\label{chap:preprocessor}

Bien que la plateforme soit une application Web (SPA), elle est adossée à un puissant pipeline de prétraitement local écrit en Python (\texttt{/preprocess}). 
L'objectif de ce pipeline est d'ingérer des fichiers microscopiques bruts (souvent volumineux et incompatibles avec les navigateurs, comme les OME-TIFF) et de les convertir en formats optimisés pour le streaming WebGL.

\section{Génération du Catalogue (\texttt{generate\_catalog.py})}
Ce script est la source de vérité de l'application Web. Il scanne le dossier racine de données et parse les métadonnées de chaque sous-dossier pour générer le \texttt{catalog.json} final consommé par la page \textbf{Explorer}.
Il extrait de manière critique :
\begin{itemize}
    \item Les dimensions $X, Y, Z, T$.
    \item La résolution physique (Voxel Size en $\mu m$) utilisée pour calibrer l'échelle de l'interface (Scale Bar) et les outils de mesure.
    \item Les métadonnées biologiques (Stade de développement, date, couleurs des canaux).
\end{itemize}

\section{Tuilage et Bricking Out-of-Core (Dask)}
Pour les volumes dépassant la mémoire vive (RAM) disponible, le pipeline utilise la bibliothèque \textbf{Dask} (\texttt{brick\_preprocessor\_dask.py}).
Plutôt que de charger un volume entier de 10 Go en mémoire, le volume est subdivisé ("chunké") en petits blocs tridimensionnels (Bricks).
Le script \texttt{pack\_bricks.py} compresse ensuite ces données dans un format binaire optimisé que le \texttt{VolumeViewer} (WebGL) peut streamer progressivement (Progressive Loading).

\section{Stabilisation 3D (Kabsch Algorithm)}
Pour les données de Tracking temporel ("Live"), les embryons ont tendance à dériver ou tourner dans le champ du microscope au fil du temps.
Le script \texttt{registration.py} implémente un algorithme de \textbf{Rigid-Body Registration} :
\begin{enumerate}
    \item Il lit le fichier \texttt{tracks.json} brut contenant les coordonnées $X, Y, Z$ de toutes les cellules pour chaque instant $t$.
    \item Il identifie des paires de points homologues (cellules de référence) entre l'instant $t$ et l'instant $t=0$.
    \item Il applique l'\textbf{Algorithme de Kabsch} via Singular Value Decomposition (SVD) :
    \begin{equation}
        \text{Covariance} = \text{Source}_{\text{centrée}}^T \cdot \text{Target}_{\text{centrée}}
    \end{equation}
    \begin{equation}
        U, S, V^T = \text{SVD}(\text{Covariance})
    \end{equation}
    \begin{equation}
        R = V^T \cdot U^T
    \end{equation}
    \item Une routine de rejet d'outliers (Median Absolute Deviation) filtre les cellules erratiques (ex: cellules mortes ou en division) pour ne conserver que le mouvement global de la masse tissulaire.
    \item Le script exporte les matrices de Rotation ($R$) et de Translation ($T$) permettant au visualiseur d'afficher un embryon parfaitement stabilisé au centre de l'écran.
\end{enumerate}

\section{Optimisation Web (\textit{DeepZoom} et Thumbnails)}
Enfin, \texttt{generate\_deepzoom\_tiles.py} construit des pyramides de résolutions 2D pour le composant \texttt{DeepZoomViewer}. Le processus génère un \texttt{manifest.json} et des dizaines de milliers de petites tuiles compressées (souvent en WEBP pour minimiser la bande passante), permettant un chargement fluide sur la page \textbf{DeepZoom}.
